%%%%%%%%%%%%%%%%%%%%%%%%%%%%%%%%%%%%%%%%%%%%%%%%%%%%%%%%%%%%%%%%%%%%%%%%%%%
% Función: capítulo de ejemplo de objetivos de la tesis por compendio.
%
% Uso: modifica este fichero para adaptar el contenido a tu trabajo.
%%%%%%%%%%%%%%%%%%%%%%%%%%%%%%%%%%%%%%%%%%%%%%%%%%%%%%%%%%%%%%%%%%%%%%%%%%%

\chapter{Objetivos}
\label{cha:compendium-objectives}

Al formular objetivos medibles, distingue la tarea de estimación ilustrada por la localización de cara y boca~\cite{Sanabria2026Landmark}, el problema de evaluación ligado a la aplicación~\cite{nieva2024proposal} y el despliegue de sensores y análisis descritos para EYEFUL~\cite{nieva2023eyeful}. Úsalos como ejemplos de distintos niveles de objetivos, no como objetivos o resultados que tu tesis ya haya alcanzado.

\section{Objetivo general}

Expón la hipótesis principal o el objetivo general de la tesis y deriva los objetivos específicos abordados por las publicaciones. Haz explícita la relación entre estos objetivos y la organización del resumen extendido.

\section{Objetivos específicos}
\subsection{Objetivos científicos}

Identifica los avances de conocimiento perseguidos por la tesis y relaciona cada uno con las preguntas de investigación.

\subsection{Objetivos técnicos y de validación}

Describe los objetivos de implementación, experimentación y validación que permiten evaluar los objetivos científicos.

%%% Local Variables:
%%% TeX-master: "../../book"
%%% End:
